% =============================================================================
% Assignments/CS502/02-computer-system-organization-assignment.tex
% Assignment 2: Computer-System Organization and Architecture
%
% Student set — NO answers. Solution key: 02-computer-system-organization-solutions.tex.
% Sourced from Slides/CS502/02-computer-system-organization.tex; every problem
% reuses the lecture's own notation and extends its worked examples to new
% concrete instances rather than repeating them verbatim.
%
% Compile with: cd Assignments/CS502 &&
%   TEXINPUTS=../../Preambles;$TEXINPUTS xelatex -interaction=nonstopmode 02-computer-system-organization-assignment.tex
%   (x2; Windows/MiKTeX: use ; not : as the separator)
% =============================================================================

\documentclass[11pt]{article}
\usepackage[margin=1in]{geometry}
\input{../../Preambles/header}
\coursecode{CS502}
\renewcommand{\thesection}{2.\arabic{section}}

\title{Assignment 2: Computer-System Organization and Architecture}
\author{Auqib Hamid Lone\\[0.3em]
  \emph{Department of Computer Science \& Engineering}, \emph{GCET Ganderbal}}
\date{\today}

\begin{document}
\maketitle

\noindent\textbf{Deadline:} \key{two weeks from the date of issue} --- there is no automatic late penalty; if you need more time, ask before the deadline and an extension will be granted (see the course policy in the syllabus).

\noindent Answer every question. Show all working for Numerical and Design questions --- a bare final answer without the derivation earns partial credit only. State the units of every quantity you compute.

\section{Conceptual}

\textbf{Q1.} Both interrupts and traps transfer control to the kernel, so students often treat them as the same thing. They are not. Consider these four concrete events:
\begin{enumerate}
  \item[(i)] a user presses a key on the keyboard;
  \item[(ii)] a disk controller finishes the block transfer the kernel asked it to perform;
  \item[(iii)] a program executes a division instruction with a zero divisor;
  \item[(iv)] a program executes the instruction that enters the kernel to open a file.
\end{enumerate}
(a) Classify each event as an \emph{interrupt} or a \emph{trap}, and justify the classification in one sentence each. (b) State, for each event, whether it is tied to the instruction the CPU is executing or can arrive between any two instructions. (c) Explain why this distinction matters for kernel design: what can an interrupt handler assume that a trap handler cannot (or vice versa), and why does that change how each handler is written?

\textbf{Q2.} A high-throughput disk array delivers data in $4\ \text{KB}$ blocks at $400\ \text{MB/s}$, and servicing one interrupt (entry, handler, return) costs $12\ \mu\text{s}$ of CPU time. Explain why an interrupt-per-block design cannot sustain this throughput, using a per-event cost argument: state how the CPU overhead depends on the event rate and the per-event cost, give a one-line arithmetic estimate for this device, and state the general condition under which the machine saturates. Then name the mechanism that removes this per-event cost, and explain in one sentence what it does differently.

\section{Numerical}

\textbf{Q3.} A network interface raises an interrupt for every packet it receives. On a certain system, interrupt entry and exit cost $5\ \mu\text{s}$ together, and the handler itself costs $15\ \mu\text{s}$, so each interrupt consumes $20\ \mu\text{s}$ of CPU time. The interface currently delivers $6{,}000$ packets per second.
\begin{enumerate}
  \item[(a)] Compute the fraction of CPU time spent handling these interrupts, as a percentage.
  \item[(b)] At what packet rate would interrupt handling consume 50\% of the CPU?
  \item[(c)] At what packet rate would the CPU be fully saturated (100\% overhead, no useful work left)?
\end{enumerate}
Show the arithmetic and give units for each rate.

\textbf{Q4.} A magnetic disk spins at $15{,}000\ \text{RPM}$. Its average seek time is $4.5\ \text{ms}$, its sustained transfer rate is $200\ \text{MB/s}$, and the kernel reads in $64\ \text{KB}$ blocks. Use $1\ \text{KB} = 1024\ \text{B}$ and $1\ \text{MB/s} = 10^6\ \text{B/s}$.
\begin{enumerate}
  \item[(a)] Compute one full rotation, the average rotational latency, and the transfer time for one block.
  \item[(b)] Compute the total access time for one block, and the percentage of that total spent on mechanical movement (seek plus rotational latency).
  \item[(c)] Recompute the total access time if the average seek time is improved to $2.0\ \text{ms}$, and report the percentage reduction. What does this say about which component dominates?
\end{enumerate}

\section{Design}

\textbf{Q5.} You are choosing an I/O technique for each of three devices on the same machine. For each, choose \emph{polling}, \emph{interrupt-driven I/O}, or \emph{DMA}, and justify the choice with a cost argument (state the relevant event or block rate and the per-event cost, and compare it with the alternative).
\begin{enumerate}
  \item[(a)] A keyboard that produces about $15$ keystrokes per second; an interrupt costs $20\ \mu\text{s}$.
  \item[(b)] A temperature sensor that is always ready and must be sampled every $500\ \mu\text{s}$; an interrupt costs $20\ \mu\text{s}$, while one status-register poll plus the data read costs about $0.2\ \mu\text{s}$.
  \item[(c)] A disk that reads a $1\ \text{GB}$ file at $500\ \text{MB/s}$; an interrupt costs $20\ \mu\text{s}$, and DMA setup plus one completion interrupt also costs $20\ \mu\text{s}$. Consider the interrupt-per-$4\ \text{KB}$-block option and at least one larger block size.
\end{enumerate}
For (c) you may choose the block size, but state it and show the interrupt count.

\textbf{Q6.} Design the I/O path for a 4-core machine reading a large file from a disk while all four cores run user code. The machine has \emph{one} shared memory bus, and every core has private L1/L2 caches plus a shared last-level cache. Assume a $2\ \text{GB}$ file, a disk that sustains $400\ \text{MB/s}$, and a per-event CPU cost of $20\ \mu\text{s}$ (interrupt-driven) or $20\ \mu\text{s}$ per DMA setup-plus-completion.
\begin{enumerate}
  \item[(a)] Choose a block size and justify it quantitatively: compare the interrupt count and the CPU overhead for $4\ \text{KB}$ blocks against your chosen size over the whole file.
  \item[(b)] Choose DMA or interrupt-driven I/O, with the arithmetic to support the choice.
  \item[(c)] Name one cache-coherence risk this design creates, explain the concrete sequence of events that produces it, and state how the OS and hardware prevent it.
  \item[(d)] State the bottleneck the single memory bus creates, and one design change that would relieve it.
\end{enumerate}

\end{document}
